% EDITED 2026-10-18 scripts/patch_continuity_appendix_linebreaks_2026-10-18T2130.py
% EDITED 2026-10-18 scripts/patch_continuity_review_appendix_2026-10-18T2119.py
% BEGIN CONTINUITY SELF REFERENCE APPENDIX
% Continuity of reversible transformations from self-reference.
% Adapted from knowledge/sessions/
% continuity_from_self_reference_open_problem_2026-10-18T1811.md.
% Status labels are proposed for JS's review; no ledger is amended.
% Requires the packages and proposition/proof environments of main.tex.
% Integration, within the existing appendices environment:
%   % EDITED 2026-10-18 scripts/patch_continuity_appendix_linebreaks_2026-10-18T2130.py
% EDITED 2026-10-18 scripts/patch_continuity_review_appendix_2026-10-18T2119.py
% BEGIN CONTINUITY SELF REFERENCE APPENDIX
% Continuity of reversible transformations from self-reference.
% Adapted from knowledge/sessions/
% continuity_from_self_reference_open_problem_2026-10-18T1811.md.
% Status labels are proposed for JS's review; no ledger is amended.
% Requires the packages and proposition/proof environments of main.tex.
% Integration, within the existing appendices environment:
%   % EDITED 2026-10-18 scripts/patch_continuity_appendix_linebreaks_2026-10-18T2130.py
% EDITED 2026-10-18 scripts/patch_continuity_review_appendix_2026-10-18T2119.py
% BEGIN CONTINUITY SELF REFERENCE APPENDIX
% Continuity of reversible transformations from self-reference.
% Adapted from knowledge/sessions/
% continuity_from_self_reference_open_problem_2026-10-18T1811.md.
% Status labels are proposed for JS's review; no ledger is amended.
% Requires the packages and proposition/proof environments of main.tex.
% Integration, within the existing appendices environment:
%   % EDITED 2026-10-18 scripts/patch_continuity_appendix_linebreaks_2026-10-18T2130.py
% EDITED 2026-10-18 scripts/patch_continuity_review_appendix_2026-10-18T2119.py
% BEGIN CONTINUITY SELF REFERENCE APPENDIX
% Continuity of reversible transformations from self-reference.
% Adapted from knowledge/sessions/
% continuity_from_self_reference_open_problem_2026-10-18T1811.md.
% Status labels are proposed for JS's review; no ledger is amended.
% Requires the packages and proposition/proof environments of main.tex.
% Integration, within the existing appendices environment:
%   \input{appendices/appendix_continuity_self_reference_2026-10-18T2119}
% Main-text reference: Appendix~\ref{app:continuity-self-reference}.
% Five citation keys await import of
% bib/continuity_open_problem_refs_2026-10-18T1811.bib and BBT re-export:
% backens2014zx, bennett2003notes, simon2001nosignaling,
% chiribella2010purification, chiribella2011informational.
% The two Chiribella et al. citations support only the limited statements
% verified from their abstracts; no detailed reconstruction audit is claimed.
%
% SUGGESTED MAIN-TEXT SUMMARY (for the later main-text edit):
% Continuous reversible accessibility remains \textsc{open}. Complex
% amplitudes, probability normalisation and symmetry among admitted bases
% do not individually establish physically available paths between pure
% states. With a closed, transitive reversible group on a finite-dimensional
% compact convex state body, connectedness of the pure states would suffice;
% the outstanding task is to justify that connectedness physically.
%
% \textsc{Candidate.} The missing link may concern the complete act of
% self-description: producing a record changes the observer, and the memory,
% its encoding and the subsequent feedback must belong to the same physical
% account. A constraint on information, disturbance and recoverable
% correlations might connect the failure of complete self-recording to the
% alternatives the theory must admit. That connection is not yet a theorem.
% Appendix~\ref{app:continuity-self-reference} reviews the attempted routes,
% retains their valid conditional results, and gives explicit criteria for
% either a derivation or a countermodel. The argument must precede the Born
% rule and the physical $L^2$ metric.

\section{Continuity from self-reference: attempted derivations and the open problem}
\label{app:continuity-self-reference}

\noindent\textsc{Open.}
Can the physical limitations on an embedded observer's self-description
explain why any two pure states can be connected by a continuous family of
physically available reversible transformations? This appendix reviews the
attempts underlying Open Problem~\ref{op:limit-realisability}. Some proposed
implications fail; others give useful conditional mathematics while leaving
a physical premise unproved. The distinction matters: replacing that
premise by a mathematical reformulation does not explain why nature must
satisfy it. A status label applies to the discussion it introduces until
a different status is stated.

\subsection{The target and the order of dependence}
\label{app:continuity-target}

\noindent\textsc{Structural.}
Hardy's fifth axiom asks for continuous reversible accessibility between
any two pure states \autocite{Hardy2001a}. Let $P$ denote the admitted pure
states and $G$ the physically allowed reversible transformations. The target is
\begin{equation}
  \forall x,y\in P\quad\exists\,g:[0,1]\longrightarrow G
  \quad\text{continuous},\qquad g(0)=e,\quad g(1)x=y.
  \label{eq:continuity-accessibility}
\end{equation}
This is stronger than continuity of an individual map on states: the
intermediate transformations must also be available. The parameter need
not already be a physical time coordinate, and differentiability is not
an initial premise.

A topology can be specified before the Born rule. In a finite-dimensional
operational description, a finite list of measurement probabilities
separates states, and convergence means convergence of those probabilities.
Pure states are the extreme points of the admitted convex preparation set,
not automatically all rays of a vector space. This specifies what
connectedness would mean without assuming it \autocite{Hardy2001a}.

\noindent\textsc{Candidate.}
The earlier complex-amplitude construction is taken here as a conditional
starting point. The proposed order in \S\ref{sec:reversible-continuity-l2} is
\begin{equation}
\begin{aligned}
 &\text{self-reference and complex amplitudes}\\
 &\qquad\longrightarrow\text{continuous reversible accessibility}\\
 &\qquad\longrightarrow L^2
   \longrightarrow (S^1\longrightarrow S^3\longrightarrow S^2)
   \longrightarrow\text{Born probabilities}.
\end{aligned}
\label{eq:continuity-dependency}
\end{equation}
Each downstream arrow retains its own premises, including assumptions
about how states are represented and which transformations preserve the norm. Neither the
Born rule nor a physically normalised Hilbert sphere can therefore justify
the earlier continuity step.

\subsection{What self-recording prevents}
\label{app:continuity-self-recording}

\noindent\textsc{Structural.}
Writing a new record into an observer's memory changes its total physical
state, schematically $s\xrightarrow{\mathrm{write}\ r}s_r$. The record need
not thereby become false: it can describe a past state, a partial property,
or a predictable consequence of the writing. The stronger obstruction
concerns a universally successful predictor whose predictions can themselves
enter the behaviour predicted. Popper includes calculation and the physical
production of predictions in his analysis and establishes limitations in a
classical setting \autocite{Popper1950a}. This obstruction alone therefore
does not distinguish quantum theory from classical theories.

\noindent\textsc{Candidate.}
\textsc{Fanout} denotes the formation of an accessible, copyable outcome
record. The proposed pre-\textsc{fanout} regime lacks a committed outcome
record for the interaction in question. Turning that statement into the
absence of every preferred basis or invariant distinction requires an
additional argument. In particular, the impossibility of complete
self-recording is not the impossibility of recording any property.

\noindent\textsc{Rhyme.}
The liar and truth-teller suggest a comparison between a Boolean
self-assignment that is inconsistent and one that is consistent but
underdetermined. This
motivates the description of pre-\textsc{fanout} paradox and post-\textsc{fanout}
hypodox. It supplies neither a topology nor a probability law for the record
selected. The comparison with information in a closed timelike loop has
the same limited role.

\noindent\textsc{Structural.}
Inability to predict an outcome from the observer's information does not
establish absence of physical antecedents. A restriction on prediction
cannot alone support that causal conclusion or a continuity theorem.

\subsection{Attempted routes and their remaining gaps}
\label{app:continuity-routes}

\subsubsection{Complex amplitudes provide a continuum}
\label{app:continuity-complex}

\noindent\textsc{Structural.}
A continuous scalar field permits continuous families of vectors and
matrices; it does not make all of them physical. The complex stabiliser
fragment supplies a direct counterexample to that inference. Its
single-qubit pure states are the six rays represented by
$(1,0)$, $(0,1)$, $(1,\pm1)$ and $(1,\pm i)$, while its Clifford
transformations have 24 distinct actions on rays. It has complex amplitudes
and reversible linear maps without a continuous path between distinct pure
states \autocite{backens2014zx}.

Spekkens's toy bit is an earlier finite example, but is not itself a native
complex-amplitude model \autocite{spekkens2007toy}; the stabiliser fragment
addresses that objection. These examples test particular implications.
They have not thereby been checked against a fully specified physical model
of a self-describing observer.

\subsubsection{Probability normalisation forces unitarity and continuity}
\label{app:continuity-normalisation}

\noindent\textsc{Structural.}
For an exhaustive measurement, normalisation gives
\begin{equation}
  \sum_i p(i\mid s,M)=1.
  \label{eq:continuity-operational-normalisation}
\end{equation}
These are real outcome probabilities. Identifying them with $|a_i|^2$
already supplies the quadratic rule that the proposed order of dependence
places later. The attempted conservation argument made this identification
as a premise; its calculation is conditional on that premise and cannot
establish the earlier continuity step.

Even granting a quadratic norm and complex-linear dynamics, norm
preservation gives $U^\dagger U=I$, not the physical availability of every
element of $U(n)$. A finite subgroup also preserves the norm. Continuity of
each linear map remains distinct from a continuous path of allowed maps
connecting it to the identity. Both gaps survive the normalisation calculation.

\subsubsection{No record means no preferred basis and smooth access between bases}
\label{app:continuity-bases}

\noindent\textsc{Structural.}
Absence of a particular record does not exclude a restriction in the laws.
Symmetry among admitted bases does not imply that every mathematical basis
is admitted. Transitivity among admitted states does not imply continuous
transitivity: a finite group can move any member of a finite set to any
other. A passive change of coordinates is also distinct from a physical
operation. The stabiliser fragment treats its three Pauli bases symmetrically
while restricting the family of bases \autocite{backens2014zx}.

An invariant mathematical partition is not automatically a readable
classical record. An argument using that identification must construct an
allowed measurement that distinguishes the parts and an allowed procedure
that copies the result. Naming the parts supplies neither operation. Thus
the claim that continuity forbids a privileged finite basis does not
establish its converse.

\subsubsection{No thermodynamic cost means reversible evolution}
\label{app:continuity-thermodynamics}

\noindent\textsc{Structural.}
Logical invertibility, absence of heat dissipation and thermodynamic
reversibility are different conditions. Landauer reasoning tracks
information transferred to other degrees of freedom; it does not derive
microscopic reversible laws from zero heat. Copying classical data into a
blank register can itself be implemented reversibly \autocite{bennett2003notes}.
A cost assigned to \textsc{fanout} consequently cannot by itself prove
that all pre-\textsc{fanout} operations possess physically available
inverses, still less that those inverses belong to continuous paths.

\subsubsection{Excluded superpositions cause nonlinear dynamics and signalling}
\label{app:continuity-signalling}

\noindent\textsc{Structural.}
The first implication is false. The complex-linear reversible matrix
$S=\operatorname{diag}(1,i)$ preserves the six stabiliser rays; restricting
the state set does not make that matrix nonlinear \autocite{backens2014zx}.
A restriction to standard quantum preparations, operations and measurements
cannot create a signalling protocol unavailable in the full theory. These
known quantum facts provide a counterexample to the inference, not premises
for the reconstruction.

Simon, Bužek and Gisin derive linear, completely positive density-matrix
dynamics using Hilbert-space states, the usual trace probability rule, and
no signalling \autocite{simon2001nosignaling}. Their theorem cannot be moved
upstream of the metric and Born rule without replacing those premises.
An independently formulated operational no-signalling condition does not
by itself exclude the restricted subtheory just described.

\subsubsection{Diagonalisation or inexhaustible self-ignorance supplies continuity}
\label{app:continuity-diagonal}

\noindent\textsc{Structural.}
Diagonalisation establishes failure of exhaustive representation under
specified assumptions, not connectedness \autocite{Yanofsky2003}. For example,
the space of infinite binary sequences is uncountable but totally disconnected:
two distinct sequences differ at a coordinate whose two values separate
them into disjoint open-and-closed sets. Uncountably many possibilities
need not admit a nonconstant continuous path. Density in a larger space
also does not establish that missing limits are physically admitted.

A continuum of outcome probabilities is insufficient for a different
reason. A classical bit admits every $p\in[0,1]$ through mixtures while
retaining two pure states. Likewise, analytic continuation of a logarithm
requires a specified domain and regularity; it does not establish that
physical amplitudes occupy that domain. The value-continuum premise used
in the Shannon/logarithm construction therefore cannot also count as a
derivation of connected pure states.

\subsubsection{Observer-dependent pure and mixed descriptions provide purification}
\label{app:continuity-purification}

\noindent\textsc{Structural.}
Conditioning must be distinguished from taking a subsystem. After Alice
measures her half of $(|00\rangle+|11\rangle)/\sqrt{2}$ in the computational
basis and learns the result, she assigns Bob the corresponding conditional
pure state. Without the result, Bob assigns his qubit $I/2$. A pure
description of the combined system does not make Bob's local state pure.
These standard quantum facts clarify the example; they are not pre-Born
premises for the continuity argument.

Hilbert-space purification says that every mixed state is the marginal of a
pure state of a larger system; its usual construction already uses the
Hilbert structure. Operational purification can instead be postulated before
that structure, as in Chiribella, D'Ariano and Perinotti's programme
\autocite{chiribella2010purification}. Their reconstruction also uses other
informational axioms \autocite{chiribella2011informational}. Purification has
not here been derived from the impossibility of complete self-description.

\subsubsection{No distinguished step implies arbitrary physical roots}
\label{app:continuity-roots}

\noindent\textsc{Structural.}
Under the finite-dimensional compact-group assumptions below, divisibility
of $G$ --- every element has an $n$-th root for every positive integer $n$ ---
implies connectedness. The gap is physical: absence of a distinguished step
does not establish the availability of all those roots. A formal matrix
root need not be an allowed operation. Proposition~\ref{prop:continuity-review-divisible}
retains the conditional result without supplying that premise.

\subsubsection{Connected pure states and symmetry imply continuous reversibility}
\label{app:continuity-connected}

\noindent\textsc{Structural.}
This route works conditionally. A closed group of reversible affine
transformations of a finite-dimensional compact convex state body, acting
transitively on connected pure states, has an identity component that
already acts transitively. Proposition~\ref{prop:continuity-review-connected}
gives the proof. The physical origin of connectedness remains unproved,
with closedness and transitivity also requiring explicit treatment.

Admitting all complex directions would provide a connected space even
before choosing $L^2$: the unit sphere of any norm on $\mathbb C^n$, for
$n\geq1$, is connected by radial rescaling of an ordinary sphere. A complex
representation, however, does not establish that admission. Closure under
limits does not suffice either: finite state sets and finite transformation
groups are already closed.

\subsubsection{Self-consistency requires pure fixed points}
\label{app:continuity-fixed-points}

\noindent\textsc{Structural.}
The fixed-point route proposes
\begin{equation}
  \forall g\in G\quad\exists x\in P:\ gx=x.
  \label{eq:continuity-pure-fixed-point}
\end{equation}
With compact Lie structure, transitivity and at least two pure states,
this forces a connected, positive-dimensional pure-state space
(Proposition~\ref{prop:continuity-review-fixed-points}). It is a useful
reduction of the target.

\noindent\textsc{Open.}
What remains is the derivation of \emph{pure} fixed points from physical
self-reference. A classical bit-flip has a consistent equal mixture without
a pure fixed point. Lawvere's theorem requires a suitable representation
and evaluation structure and a surjectivity hypothesis
\autocite{Lawvere_1969,Yanofsky2003}. Calling a situation self-referential
does not verify these hypotheses or put the resulting fixed point in $P$.
This route passes only if the physical and mathematical requirements are
established without assuming the full Hilbert space. It fails if consistency
supplies only mixed states or if the representation hypothesis is merely
postulated.

\subsection{The possible missing physical insight}
\label{app:continuity-candidate}

\noindent\textsc{Candidate.}
The common gap is a relation between the information an embedded observer
cannot record and the physical alternatives a theory must admit. An
obstruction to a universally correct record says what cannot be completed;
continuous reversibility asserts the existence of a family of processes.
A physical principle must explain that existence.

The complete act of self-description is a possible place to seek it. The
target, the memory describing it, the physical encoding and the feedback
that uses the description must belong to the same account. Adding memory
changes the target; redescribing the enlarged target is another physical
act. The relevant constraint may concern information, disturbance and
recoverable correlations among these parts when complete internal recording
is unavailable. It has not been shown that these correlations are phases,
that they form a continuum, or that their transformations are reversible.

One possible intermediate result would link an obstruction to continuous
reversible access to an internally available, repeatably readable distinction.
If that link followed from the observer model, the absence of the relevant
record could constrain the decomposition of pure states. The measurement
and copying operations must be derived; a partition alone does not supply
them. Another possibility is that the same analysis forces an operational
purification or reversible-dilation principle. That direction would require
an audit of the remaining reconstruction premises. Neither possibility is
presently a theorem.

\subsection{The open problem in operational form}
\label{app:continuity-open}

\noindent\textsc{Open.}
The following formulation sharpens Open Problem~\ref{op:limit-realisability}.
Work with a finite-dimensional operational theory. Declare that preparations
form a compact convex set $\Omega$, a finite collection of measurement
probabilities separates its states, and deterministic reversible operations
act affinely. Put $P=\operatorname{ext}(\Omega)$. Take complex amplitudes
only with the algebraic and compositional properties actually obtained
from the earlier self-reference argument; do not add full superposition
admissibility.

\paragraph{Physical premise.}
Specify an operational condition $\mathcal R$ that states the limits of
an observer's self-description. Include the observer, memory, record production and
feedback in the same theory. Identify the prediction tasks admitting a
diagonal response and the partial or retrospective records that remain
possible. Distinguish closure of the physical model under adding observers
from an observer's possession of an exhaustive internal representation.
Where more memory is needed, use a family of finite systems rather than a
finite register endowed with unlimited capacity.

\paragraph{Narrow mathematical task.}
Initially grant that $G$ is closed and transitive on $P$, explicitly as
additional premises. Determine whether $\mathcal R$, together with the
declared complex-amplitude structure, forces $P$ to be connected, or
equivalently under these assumptions forces $G_0$ to act transitively.
Deriving \eqref{eq:continuity-pure-fixed-point} would be one sufficient
route, but is not required.

\paragraph{Full physical task.}
Explain whether closedness, transitivity and available inverses also
follow from the observer analysis. Identify any that remain independent
postulates. A solution of the narrow task must not be presented as deriving
continuous reversibility from self-reference alone.

\paragraph{PASS: affirmative resolution.}
State $\mathcal R$ in independently meaningful operational terms and prove
\eqref{eq:continuity-accessibility} for every model of the declared premises.
Establish the available paths or sufficient connectedness result and identify
the precise premise violated by the finite counterexamples. Constructing
one connected model does not establish necessity.

\paragraph{PASS: negative resolution.}
Construct a model satisfying the same explicit premises, including the
physical observer and its prediction and record procedures, in which the
target fails. For the narrow task, retain closedness and transitivity. The
model must satisfy the substantive complex-amplitude requirements, not
merely encode classical labels with complex coordinates. Such a model would
establish that an additional physical principle is necessary.

\paragraph{FAIL as a derivation.}
Import the Born rule, a physical $L^2$ sphere, all projective rays, all
unitary transformations, continuous control or arbitrary physical roots as
an unstated premise; infer pure-state connectedness merely from continuously
variable probabilities; or assert that a toy model violates $\mathcal R$
without defining and testing it. A premise that only rephrases connectedness
also fails the explanatory aim even if the subsequent mathematics is correct.

\subsection{Conditional mathematical results}
\label{app:continuity-conditional-results}

\noindent\textsc{Structural.}
Let $\Omega$ be a compact convex set spanning a finite-dimensional affine
space and $G$ a closed subgroup of its affine automorphisms. The
automorphism group is compact: its centroid is a common fixed point, and
inscribed and circumscribed balls bound the matrices and their inverses.
The closed-subgroup theorem makes $G$ a compact matrix Lie group; its
identity component $G_0$ is path-connected and has finite index. Auxiliary
Euclidean coordinates in this argument do not identify physical
normalisation or probabilities with a quadratic norm. All three propositions
below use these hypotheses; arbitrary compact groups need not be Lie groups.

\begin{proposition}[Connected pure states]
\label{prop:continuity-review-connected}
If $G$ acts transitively on a connected pure-state space $P$, then $G_0$
acts transitively on $P$ and \eqref{eq:continuity-accessibility} holds.
\end{proposition}
\begin{proof}
The $G_0$-orbits partition $P$ into finitely many sets because $G/G_0$ is
finite and $G$ is transitive. Each orbit is compact and hence closed; in
this finite partition each is also open. Connectedness leaves one orbit.
For $x,y\in P$, choose $h\in G_0$ with $hx=y$ and a continuous path in
$G_0$ from $e$ to $h$. This is the required family of transformations.
The whole group $G$ need not be connected.
\end{proof}

\begin{proposition}[Divisibility]
\label{prop:continuity-review-divisible}
If $G$ is divisible, it is connected. If it also acts transitively on $P$,
then \eqref{eq:continuity-accessibility} holds.
\end{proposition}
\begin{proof}
The finite quotient $G/G_0$ inherits divisibility. In a finite group of
order $m$, every element satisfies $h^m=e$, so a nonidentity element has
no $m$-th root. The quotient must therefore be trivial. A connected Lie
group is path-connected, and transitivity gives the result.
\end{proof}

\begin{proposition}[Pure fixed points]
\label{prop:continuity-review-fixed-points}
Suppose $G$ acts transitively on $P$, with $|P|\geq2$, and every $g\in G$
fixes some point of $P$. Then $P$ is a connected manifold of positive
dimension and \eqref{eq:continuity-accessibility} holds.
\end{proposition}
\begin{proof}
The finite permutation group induced on the $G_0$-orbits is transitive.
If there were more than one orbit, this group would have an element moving
every orbit: its average number of fixed points is one by the orbit-counting
lemma, while its identity fixes more than one. Any lift to $G$ would fix no
point of $P$, contradicting the hypothesis. Thus $P$ is a single $G_0$-orbit
and is connected. Being a nontrivial homogeneous space of a compact Lie
group, it has positive dimension. Proposition~\ref{prop:continuity-review-connected}
supplies the continuous reversible paths.
\end{proof}

Averaging any state over a compact group gives an invariant state of
$\Omega$, which may be mixed. Consistency in the convex state space is
therefore much weaker than the pure fixed-point premise. A physical reason
to require an admissible pure fixed point is precisely what that route
still owes.
% END CONTINUITY SELF REFERENCE APPENDIX

% Main-text reference: Appendix~\ref{app:continuity-self-reference}.
% Five citation keys await import of
% bib/continuity_open_problem_refs_2026-10-18T1811.bib and BBT re-export:
% backens2014zx, bennett2003notes, simon2001nosignaling,
% chiribella2010purification, chiribella2011informational.
% The two Chiribella et al. citations support only the limited statements
% verified from their abstracts; no detailed reconstruction audit is claimed.
%
% SUGGESTED MAIN-TEXT SUMMARY (for the later main-text edit):
% Continuous reversible accessibility remains \textsc{open}. Complex
% amplitudes, probability normalisation and symmetry among admitted bases
% do not individually establish physically available paths between pure
% states. With a closed, transitive reversible group on a finite-dimensional
% compact convex state body, connectedness of the pure states would suffice;
% the outstanding task is to justify that connectedness physically.
%
% \textsc{Candidate.} The missing link may concern the complete act of
% self-description: producing a record changes the observer, and the memory,
% its encoding and the subsequent feedback must belong to the same physical
% account. A constraint on information, disturbance and recoverable
% correlations might connect the failure of complete self-recording to the
% alternatives the theory must admit. That connection is not yet a theorem.
% Appendix~\ref{app:continuity-self-reference} reviews the attempted routes,
% retains their valid conditional results, and gives explicit criteria for
% either a derivation or a countermodel. The argument must precede the Born
% rule and the physical $L^2$ metric.

\section{Continuity from self-reference: attempted derivations and the open problem}
\label{app:continuity-self-reference}

\noindent\textsc{Open.}
Can the physical limitations on an embedded observer's self-description
explain why any two pure states can be connected by a continuous family of
physically available reversible transformations? This appendix reviews the
attempts underlying Open Problem~\ref{op:limit-realisability}. Some proposed
implications fail; others give useful conditional mathematics while leaving
a physical premise unproved. The distinction matters: replacing that
premise by a mathematical reformulation does not explain why nature must
satisfy it. A status label applies to the discussion it introduces until
a different status is stated.

\subsection{The target and the order of dependence}
\label{app:continuity-target}

\noindent\textsc{Structural.}
Hardy's fifth axiom asks for continuous reversible accessibility between
any two pure states \autocite{Hardy2001a}. Let $P$ denote the admitted pure
states and $G$ the physically allowed reversible transformations. The target is
\begin{equation}
  \forall x,y\in P\quad\exists\,g:[0,1]\longrightarrow G
  \quad\text{continuous},\qquad g(0)=e,\quad g(1)x=y.
  \label{eq:continuity-accessibility}
\end{equation}
This is stronger than continuity of an individual map on states: the
intermediate transformations must also be available. The parameter need
not already be a physical time coordinate, and differentiability is not
an initial premise.

A topology can be specified before the Born rule. In a finite-dimensional
operational description, a finite list of measurement probabilities
separates states, and convergence means convergence of those probabilities.
Pure states are the extreme points of the admitted convex preparation set,
not automatically all rays of a vector space. This specifies what
connectedness would mean without assuming it \autocite{Hardy2001a}.

\noindent\textsc{Candidate.}
The earlier complex-amplitude construction is taken here as a conditional
starting point. The proposed order in \S\ref{sec:reversible-continuity-l2} is
\begin{equation}
\begin{aligned}
 &\text{self-reference and complex amplitudes}\\
 &\qquad\longrightarrow\text{continuous reversible accessibility}\\
 &\qquad\longrightarrow L^2
   \longrightarrow (S^1\longrightarrow S^3\longrightarrow S^2)
   \longrightarrow\text{Born probabilities}.
\end{aligned}
\label{eq:continuity-dependency}
\end{equation}
Each downstream arrow retains its own premises, including assumptions
about how states are represented and which transformations preserve the norm. Neither the
Born rule nor a physically normalised Hilbert sphere can therefore justify
the earlier continuity step.

\subsection{What self-recording prevents}
\label{app:continuity-self-recording}

\noindent\textsc{Structural.}
Writing a new record into an observer's memory changes its total physical
state, schematically $s\xrightarrow{\mathrm{write}\ r}s_r$. The record need
not thereby become false: it can describe a past state, a partial property,
or a predictable consequence of the writing. The stronger obstruction
concerns a universally successful predictor whose predictions can themselves
enter the behaviour predicted. Popper includes calculation and the physical
production of predictions in his analysis and establishes limitations in a
classical setting \autocite{Popper1950a}. This obstruction alone therefore
does not distinguish quantum theory from classical theories.

\noindent\textsc{Candidate.}
\textsc{Fanout} denotes the formation of an accessible, copyable outcome
record. The proposed pre-\textsc{fanout} regime lacks a committed outcome
record for the interaction in question. Turning that statement into the
absence of every preferred basis or invariant distinction requires an
additional argument. In particular, the impossibility of complete
self-recording is not the impossibility of recording any property.

\noindent\textsc{Rhyme.}
The liar and truth-teller suggest a comparison between a Boolean
self-assignment that is inconsistent and one that is consistent but
underdetermined. This
motivates the description of pre-\textsc{fanout} paradox and post-\textsc{fanout}
hypodox. It supplies neither a topology nor a probability law for the record
selected. The comparison with information in a closed timelike loop has
the same limited role.

\noindent\textsc{Structural.}
Inability to predict an outcome from the observer's information does not
establish absence of physical antecedents. A restriction on prediction
cannot alone support that causal conclusion or a continuity theorem.

\subsection{Attempted routes and their remaining gaps}
\label{app:continuity-routes}

\subsubsection{Complex amplitudes provide a continuum}
\label{app:continuity-complex}

\noindent\textsc{Structural.}
A continuous scalar field permits continuous families of vectors and
matrices; it does not make all of them physical. The complex stabiliser
fragment supplies a direct counterexample to that inference. Its
single-qubit pure states are the six rays represented by
$(1,0)$, $(0,1)$, $(1,\pm1)$ and $(1,\pm i)$, while its Clifford
transformations have 24 distinct actions on rays. It has complex amplitudes
and reversible linear maps without a continuous path between distinct pure
states \autocite{backens2014zx}.

Spekkens's toy bit is an earlier finite example, but is not itself a native
complex-amplitude model \autocite{spekkens2007toy}; the stabiliser fragment
addresses that objection. These examples test particular implications.
They have not thereby been checked against a fully specified physical model
of a self-describing observer.

\subsubsection{Probability normalisation forces unitarity and continuity}
\label{app:continuity-normalisation}

\noindent\textsc{Structural.}
For an exhaustive measurement, normalisation gives
\begin{equation}
  \sum_i p(i\mid s,M)=1.
  \label{eq:continuity-operational-normalisation}
\end{equation}
These are real outcome probabilities. Identifying them with $|a_i|^2$
already supplies the quadratic rule that the proposed order of dependence
places later. The attempted conservation argument made this identification
as a premise; its calculation is conditional on that premise and cannot
establish the earlier continuity step.

Even granting a quadratic norm and complex-linear dynamics, norm
preservation gives $U^\dagger U=I$, not the physical availability of every
element of $U(n)$. A finite subgroup also preserves the norm. Continuity of
each linear map remains distinct from a continuous path of allowed maps
connecting it to the identity. Both gaps survive the normalisation calculation.

\subsubsection{No record means no preferred basis and smooth access between bases}
\label{app:continuity-bases}

\noindent\textsc{Structural.}
Absence of a particular record does not exclude a restriction in the laws.
Symmetry among admitted bases does not imply that every mathematical basis
is admitted. Transitivity among admitted states does not imply continuous
transitivity: a finite group can move any member of a finite set to any
other. A passive change of coordinates is also distinct from a physical
operation. The stabiliser fragment treats its three Pauli bases symmetrically
while restricting the family of bases \autocite{backens2014zx}.

An invariant mathematical partition is not automatically a readable
classical record. An argument using that identification must construct an
allowed measurement that distinguishes the parts and an allowed procedure
that copies the result. Naming the parts supplies neither operation. Thus
the claim that continuity forbids a privileged finite basis does not
establish its converse.

\subsubsection{No thermodynamic cost means reversible evolution}
\label{app:continuity-thermodynamics}

\noindent\textsc{Structural.}
Logical invertibility, absence of heat dissipation and thermodynamic
reversibility are different conditions. Landauer reasoning tracks
information transferred to other degrees of freedom; it does not derive
microscopic reversible laws from zero heat. Copying classical data into a
blank register can itself be implemented reversibly \autocite{bennett2003notes}.
A cost assigned to \textsc{fanout} consequently cannot by itself prove
that all pre-\textsc{fanout} operations possess physically available
inverses, still less that those inverses belong to continuous paths.

\subsubsection{Excluded superpositions cause nonlinear dynamics and signalling}
\label{app:continuity-signalling}

\noindent\textsc{Structural.}
The first implication is false. The complex-linear reversible matrix
$S=\operatorname{diag}(1,i)$ preserves the six stabiliser rays; restricting
the state set does not make that matrix nonlinear \autocite{backens2014zx}.
A restriction to standard quantum preparations, operations and measurements
cannot create a signalling protocol unavailable in the full theory. These
known quantum facts provide a counterexample to the inference, not premises
for the reconstruction.

Simon, Bužek and Gisin derive linear, completely positive density-matrix
dynamics using Hilbert-space states, the usual trace probability rule, and
no signalling \autocite{simon2001nosignaling}. Their theorem cannot be moved
upstream of the metric and Born rule without replacing those premises.
An independently formulated operational no-signalling condition does not
by itself exclude the restricted subtheory just described.

\subsubsection{Diagonalisation or inexhaustible self-ignorance supplies continuity}
\label{app:continuity-diagonal}

\noindent\textsc{Structural.}
Diagonalisation establishes failure of exhaustive representation under
specified assumptions, not connectedness \autocite{Yanofsky2003}. For example,
the space of infinite binary sequences is uncountable but totally disconnected:
two distinct sequences differ at a coordinate whose two values separate
them into disjoint open-and-closed sets. Uncountably many possibilities
need not admit a nonconstant continuous path. Density in a larger space
also does not establish that missing limits are physically admitted.

A continuum of outcome probabilities is insufficient for a different
reason. A classical bit admits every $p\in[0,1]$ through mixtures while
retaining two pure states. Likewise, analytic continuation of a logarithm
requires a specified domain and regularity; it does not establish that
physical amplitudes occupy that domain. The value-continuum premise used
in the Shannon/logarithm construction therefore cannot also count as a
derivation of connected pure states.

\subsubsection{Observer-dependent pure and mixed descriptions provide purification}
\label{app:continuity-purification}

\noindent\textsc{Structural.}
Conditioning must be distinguished from taking a subsystem. After Alice
measures her half of $(|00\rangle+|11\rangle)/\sqrt{2}$ in the computational
basis and learns the result, she assigns Bob the corresponding conditional
pure state. Without the result, Bob assigns his qubit $I/2$. A pure
description of the combined system does not make Bob's local state pure.
These standard quantum facts clarify the example; they are not pre-Born
premises for the continuity argument.

Hilbert-space purification says that every mixed state is the marginal of a
pure state of a larger system; its usual construction already uses the
Hilbert structure. Operational purification can instead be postulated before
that structure, as in Chiribella, D'Ariano and Perinotti's programme
\autocite{chiribella2010purification}. Their reconstruction also uses other
informational axioms \autocite{chiribella2011informational}. Purification has
not here been derived from the impossibility of complete self-description.

\subsubsection{No distinguished step implies arbitrary physical roots}
\label{app:continuity-roots}

\noindent\textsc{Structural.}
Under the finite-dimensional compact-group assumptions below, divisibility
of $G$ --- every element has an $n$-th root for every positive integer $n$ ---
implies connectedness. The gap is physical: absence of a distinguished step
does not establish the availability of all those roots. A formal matrix
root need not be an allowed operation. Proposition~\ref{prop:continuity-review-divisible}
retains the conditional result without supplying that premise.

\subsubsection{Connected pure states and symmetry imply continuous reversibility}
\label{app:continuity-connected}

\noindent\textsc{Structural.}
This route works conditionally. A closed group of reversible affine
transformations of a finite-dimensional compact convex state body, acting
transitively on connected pure states, has an identity component that
already acts transitively. Proposition~\ref{prop:continuity-review-connected}
gives the proof. The physical origin of connectedness remains unproved,
with closedness and transitivity also requiring explicit treatment.

Admitting all complex directions would provide a connected space even
before choosing $L^2$: the unit sphere of any norm on $\mathbb C^n$, for
$n\geq1$, is connected by radial rescaling of an ordinary sphere. A complex
representation, however, does not establish that admission. Closure under
limits does not suffice either: finite state sets and finite transformation
groups are already closed.

\subsubsection{Self-consistency requires pure fixed points}
\label{app:continuity-fixed-points}

\noindent\textsc{Structural.}
The fixed-point route proposes
\begin{equation}
  \forall g\in G\quad\exists x\in P:\ gx=x.
  \label{eq:continuity-pure-fixed-point}
\end{equation}
With compact Lie structure, transitivity and at least two pure states,
this forces a connected, positive-dimensional pure-state space
(Proposition~\ref{prop:continuity-review-fixed-points}). It is a useful
reduction of the target.

\noindent\textsc{Open.}
What remains is the derivation of \emph{pure} fixed points from physical
self-reference. A classical bit-flip has a consistent equal mixture without
a pure fixed point. Lawvere's theorem requires a suitable representation
and evaluation structure and a surjectivity hypothesis
\autocite{Lawvere_1969,Yanofsky2003}. Calling a situation self-referential
does not verify these hypotheses or put the resulting fixed point in $P$.
This route passes only if the physical and mathematical requirements are
established without assuming the full Hilbert space. It fails if consistency
supplies only mixed states or if the representation hypothesis is merely
postulated.

\subsection{The possible missing physical insight}
\label{app:continuity-candidate}

\noindent\textsc{Candidate.}
The common gap is a relation between the information an embedded observer
cannot record and the physical alternatives a theory must admit. An
obstruction to a universally correct record says what cannot be completed;
continuous reversibility asserts the existence of a family of processes.
A physical principle must explain that existence.

The complete act of self-description is a possible place to seek it. The
target, the memory describing it, the physical encoding and the feedback
that uses the description must belong to the same account. Adding memory
changes the target; redescribing the enlarged target is another physical
act. The relevant constraint may concern information, disturbance and
recoverable correlations among these parts when complete internal recording
is unavailable. It has not been shown that these correlations are phases,
that they form a continuum, or that their transformations are reversible.

One possible intermediate result would link an obstruction to continuous
reversible access to an internally available, repeatably readable distinction.
If that link followed from the observer model, the absence of the relevant
record could constrain the decomposition of pure states. The measurement
and copying operations must be derived; a partition alone does not supply
them. Another possibility is that the same analysis forces an operational
purification or reversible-dilation principle. That direction would require
an audit of the remaining reconstruction premises. Neither possibility is
presently a theorem.

\subsection{The open problem in operational form}
\label{app:continuity-open}

\noindent\textsc{Open.}
The following formulation sharpens Open Problem~\ref{op:limit-realisability}.
Work with a finite-dimensional operational theory. Declare that preparations
form a compact convex set $\Omega$, a finite collection of measurement
probabilities separates its states, and deterministic reversible operations
act affinely. Put $P=\operatorname{ext}(\Omega)$. Take complex amplitudes
only with the algebraic and compositional properties actually obtained
from the earlier self-reference argument; do not add full superposition
admissibility.

\paragraph{Physical premise.}
Specify an operational condition $\mathcal R$ that states the limits of
an observer's self-description. Include the observer, memory, record production and
feedback in the same theory. Identify the prediction tasks admitting a
diagonal response and the partial or retrospective records that remain
possible. Distinguish closure of the physical model under adding observers
from an observer's possession of an exhaustive internal representation.
Where more memory is needed, use a family of finite systems rather than a
finite register endowed with unlimited capacity.

\paragraph{Narrow mathematical task.}
Initially grant that $G$ is closed and transitive on $P$, explicitly as
additional premises. Determine whether $\mathcal R$, together with the
declared complex-amplitude structure, forces $P$ to be connected, or
equivalently under these assumptions forces $G_0$ to act transitively.
Deriving \eqref{eq:continuity-pure-fixed-point} would be one sufficient
route, but is not required.

\paragraph{Full physical task.}
Explain whether closedness, transitivity and available inverses also
follow from the observer analysis. Identify any that remain independent
postulates. A solution of the narrow task must not be presented as deriving
continuous reversibility from self-reference alone.

\paragraph{PASS: affirmative resolution.}
State $\mathcal R$ in independently meaningful operational terms and prove
\eqref{eq:continuity-accessibility} for every model of the declared premises.
Establish the available paths or sufficient connectedness result and identify
the precise premise violated by the finite counterexamples. Constructing
one connected model does not establish necessity.

\paragraph{PASS: negative resolution.}
Construct a model satisfying the same explicit premises, including the
physical observer and its prediction and record procedures, in which the
target fails. For the narrow task, retain closedness and transitivity. The
model must satisfy the substantive complex-amplitude requirements, not
merely encode classical labels with complex coordinates. Such a model would
establish that an additional physical principle is necessary.

\paragraph{FAIL as a derivation.}
Import the Born rule, a physical $L^2$ sphere, all projective rays, all
unitary transformations, continuous control or arbitrary physical roots as
an unstated premise; infer pure-state connectedness merely from continuously
variable probabilities; or assert that a toy model violates $\mathcal R$
without defining and testing it. A premise that only rephrases connectedness
also fails the explanatory aim even if the subsequent mathematics is correct.

\subsection{Conditional mathematical results}
\label{app:continuity-conditional-results}

\noindent\textsc{Structural.}
Let $\Omega$ be a compact convex set spanning a finite-dimensional affine
space and $G$ a closed subgroup of its affine automorphisms. The
automorphism group is compact: its centroid is a common fixed point, and
inscribed and circumscribed balls bound the matrices and their inverses.
The closed-subgroup theorem makes $G$ a compact matrix Lie group; its
identity component $G_0$ is path-connected and has finite index. Auxiliary
Euclidean coordinates in this argument do not identify physical
normalisation or probabilities with a quadratic norm. All three propositions
below use these hypotheses; arbitrary compact groups need not be Lie groups.

\begin{proposition}[Connected pure states]
\label{prop:continuity-review-connected}
If $G$ acts transitively on a connected pure-state space $P$, then $G_0$
acts transitively on $P$ and \eqref{eq:continuity-accessibility} holds.
\end{proposition}
\begin{proof}
The $G_0$-orbits partition $P$ into finitely many sets because $G/G_0$ is
finite and $G$ is transitive. Each orbit is compact and hence closed; in
this finite partition each is also open. Connectedness leaves one orbit.
For $x,y\in P$, choose $h\in G_0$ with $hx=y$ and a continuous path in
$G_0$ from $e$ to $h$. This is the required family of transformations.
The whole group $G$ need not be connected.
\end{proof}

\begin{proposition}[Divisibility]
\label{prop:continuity-review-divisible}
If $G$ is divisible, it is connected. If it also acts transitively on $P$,
then \eqref{eq:continuity-accessibility} holds.
\end{proposition}
\begin{proof}
The finite quotient $G/G_0$ inherits divisibility. In a finite group of
order $m$, every element satisfies $h^m=e$, so a nonidentity element has
no $m$-th root. The quotient must therefore be trivial. A connected Lie
group is path-connected, and transitivity gives the result.
\end{proof}

\begin{proposition}[Pure fixed points]
\label{prop:continuity-review-fixed-points}
Suppose $G$ acts transitively on $P$, with $|P|\geq2$, and every $g\in G$
fixes some point of $P$. Then $P$ is a connected manifold of positive
dimension and \eqref{eq:continuity-accessibility} holds.
\end{proposition}
\begin{proof}
The finite permutation group induced on the $G_0$-orbits is transitive.
If there were more than one orbit, this group would have an element moving
every orbit: its average number of fixed points is one by the orbit-counting
lemma, while its identity fixes more than one. Any lift to $G$ would fix no
point of $P$, contradicting the hypothesis. Thus $P$ is a single $G_0$-orbit
and is connected. Being a nontrivial homogeneous space of a compact Lie
group, it has positive dimension. Proposition~\ref{prop:continuity-review-connected}
supplies the continuous reversible paths.
\end{proof}

Averaging any state over a compact group gives an invariant state of
$\Omega$, which may be mixed. Consistency in the convex state space is
therefore much weaker than the pure fixed-point premise. A physical reason
to require an admissible pure fixed point is precisely what that route
still owes.
% END CONTINUITY SELF REFERENCE APPENDIX

% Main-text reference: Appendix~\ref{app:continuity-self-reference}.
% Five citation keys await import of
% bib/continuity_open_problem_refs_2026-10-18T1811.bib and BBT re-export:
% backens2014zx, bennett2003notes, simon2001nosignaling,
% chiribella2010purification, chiribella2011informational.
% The two Chiribella et al. citations support only the limited statements
% verified from their abstracts; no detailed reconstruction audit is claimed.
%
% SUGGESTED MAIN-TEXT SUMMARY (for the later main-text edit):
% Continuous reversible accessibility remains \textsc{open}. Complex
% amplitudes, probability normalisation and symmetry among admitted bases
% do not individually establish physically available paths between pure
% states. With a closed, transitive reversible group on a finite-dimensional
% compact convex state body, connectedness of the pure states would suffice;
% the outstanding task is to justify that connectedness physically.
%
% \textsc{Candidate.} The missing link may concern the complete act of
% self-description: producing a record changes the observer, and the memory,
% its encoding and the subsequent feedback must belong to the same physical
% account. A constraint on information, disturbance and recoverable
% correlations might connect the failure of complete self-recording to the
% alternatives the theory must admit. That connection is not yet a theorem.
% Appendix~\ref{app:continuity-self-reference} reviews the attempted routes,
% retains their valid conditional results, and gives explicit criteria for
% either a derivation or a countermodel. The argument must precede the Born
% rule and the physical $L^2$ metric.

\section{Continuity from self-reference: attempted derivations and the open problem}
\label{app:continuity-self-reference}

\noindent\textsc{Open.}
Can the physical limitations on an embedded observer's self-description
explain why any two pure states can be connected by a continuous family of
physically available reversible transformations? This appendix reviews the
attempts underlying Open Problem~\ref{op:limit-realisability}. Some proposed
implications fail; others give useful conditional mathematics while leaving
a physical premise unproved. The distinction matters: replacing that
premise by a mathematical reformulation does not explain why nature must
satisfy it. A status label applies to the discussion it introduces until
a different status is stated.

\subsection{The target and the order of dependence}
\label{app:continuity-target}

\noindent\textsc{Structural.}
Hardy's fifth axiom asks for continuous reversible accessibility between
any two pure states \autocite{Hardy2001a}. Let $P$ denote the admitted pure
states and $G$ the physically allowed reversible transformations. The target is
\begin{equation}
  \forall x,y\in P\quad\exists\,g:[0,1]\longrightarrow G
  \quad\text{continuous},\qquad g(0)=e,\quad g(1)x=y.
  \label{eq:continuity-accessibility}
\end{equation}
This is stronger than continuity of an individual map on states: the
intermediate transformations must also be available. The parameter need
not already be a physical time coordinate, and differentiability is not
an initial premise.

A topology can be specified before the Born rule. In a finite-dimensional
operational description, a finite list of measurement probabilities
separates states, and convergence means convergence of those probabilities.
Pure states are the extreme points of the admitted convex preparation set,
not automatically all rays of a vector space. This specifies what
connectedness would mean without assuming it \autocite{Hardy2001a}.

\noindent\textsc{Candidate.}
The earlier complex-amplitude construction is taken here as a conditional
starting point. The proposed order in \S\ref{sec:reversible-continuity-l2} is
\begin{equation}
\begin{aligned}
 &\text{self-reference and complex amplitudes}\\
 &\qquad\longrightarrow\text{continuous reversible accessibility}\\
 &\qquad\longrightarrow L^2
   \longrightarrow (S^1\longrightarrow S^3\longrightarrow S^2)
   \longrightarrow\text{Born probabilities}.
\end{aligned}
\label{eq:continuity-dependency}
\end{equation}
Each downstream arrow retains its own premises, including assumptions
about how states are represented and which transformations preserve the norm. Neither the
Born rule nor a physically normalised Hilbert sphere can therefore justify
the earlier continuity step.

\subsection{What self-recording prevents}
\label{app:continuity-self-recording}

\noindent\textsc{Structural.}
Writing a new record into an observer's memory changes its total physical
state, schematically $s\xrightarrow{\mathrm{write}\ r}s_r$. The record need
not thereby become false: it can describe a past state, a partial property,
or a predictable consequence of the writing. The stronger obstruction
concerns a universally successful predictor whose predictions can themselves
enter the behaviour predicted. Popper includes calculation and the physical
production of predictions in his analysis and establishes limitations in a
classical setting \autocite{Popper1950a}. This obstruction alone therefore
does not distinguish quantum theory from classical theories.

\noindent\textsc{Candidate.}
\textsc{Fanout} denotes the formation of an accessible, copyable outcome
record. The proposed pre-\textsc{fanout} regime lacks a committed outcome
record for the interaction in question. Turning that statement into the
absence of every preferred basis or invariant distinction requires an
additional argument. In particular, the impossibility of complete
self-recording is not the impossibility of recording any property.

\noindent\textsc{Rhyme.}
The liar and truth-teller suggest a comparison between a Boolean
self-assignment that is inconsistent and one that is consistent but
underdetermined. This
motivates the description of pre-\textsc{fanout} paradox and post-\textsc{fanout}
hypodox. It supplies neither a topology nor a probability law for the record
selected. The comparison with information in a closed timelike loop has
the same limited role.

\noindent\textsc{Structural.}
Inability to predict an outcome from the observer's information does not
establish absence of physical antecedents. A restriction on prediction
cannot alone support that causal conclusion or a continuity theorem.

\subsection{Attempted routes and their remaining gaps}
\label{app:continuity-routes}

\subsubsection{Complex amplitudes provide a continuum}
\label{app:continuity-complex}

\noindent\textsc{Structural.}
A continuous scalar field permits continuous families of vectors and
matrices; it does not make all of them physical. The complex stabiliser
fragment supplies a direct counterexample to that inference. Its
single-qubit pure states are the six rays represented by
$(1,0)$, $(0,1)$, $(1,\pm1)$ and $(1,\pm i)$, while its Clifford
transformations have 24 distinct actions on rays. It has complex amplitudes
and reversible linear maps without a continuous path between distinct pure
states \autocite{backens2014zx}.

Spekkens's toy bit is an earlier finite example, but is not itself a native
complex-amplitude model \autocite{spekkens2007toy}; the stabiliser fragment
addresses that objection. These examples test particular implications.
They have not thereby been checked against a fully specified physical model
of a self-describing observer.

\subsubsection{Probability normalisation forces unitarity and continuity}
\label{app:continuity-normalisation}

\noindent\textsc{Structural.}
For an exhaustive measurement, normalisation gives
\begin{equation}
  \sum_i p(i\mid s,M)=1.
  \label{eq:continuity-operational-normalisation}
\end{equation}
These are real outcome probabilities. Identifying them with $|a_i|^2$
already supplies the quadratic rule that the proposed order of dependence
places later. The attempted conservation argument made this identification
as a premise; its calculation is conditional on that premise and cannot
establish the earlier continuity step.

Even granting a quadratic norm and complex-linear dynamics, norm
preservation gives $U^\dagger U=I$, not the physical availability of every
element of $U(n)$. A finite subgroup also preserves the norm. Continuity of
each linear map remains distinct from a continuous path of allowed maps
connecting it to the identity. Both gaps survive the normalisation calculation.

\subsubsection{No record means no preferred basis and smooth access between bases}
\label{app:continuity-bases}

\noindent\textsc{Structural.}
Absence of a particular record does not exclude a restriction in the laws.
Symmetry among admitted bases does not imply that every mathematical basis
is admitted. Transitivity among admitted states does not imply continuous
transitivity: a finite group can move any member of a finite set to any
other. A passive change of coordinates is also distinct from a physical
operation. The stabiliser fragment treats its three Pauli bases symmetrically
while restricting the family of bases \autocite{backens2014zx}.

An invariant mathematical partition is not automatically a readable
classical record. An argument using that identification must construct an
allowed measurement that distinguishes the parts and an allowed procedure
that copies the result. Naming the parts supplies neither operation. Thus
the claim that continuity forbids a privileged finite basis does not
establish its converse.

\subsubsection{No thermodynamic cost means reversible evolution}
\label{app:continuity-thermodynamics}

\noindent\textsc{Structural.}
Logical invertibility, absence of heat dissipation and thermodynamic
reversibility are different conditions. Landauer reasoning tracks
information transferred to other degrees of freedom; it does not derive
microscopic reversible laws from zero heat. Copying classical data into a
blank register can itself be implemented reversibly \autocite{bennett2003notes}.
A cost assigned to \textsc{fanout} consequently cannot by itself prove
that all pre-\textsc{fanout} operations possess physically available
inverses, still less that those inverses belong to continuous paths.

\subsubsection{Excluded superpositions cause nonlinear dynamics and signalling}
\label{app:continuity-signalling}

\noindent\textsc{Structural.}
The first implication is false. The complex-linear reversible matrix
$S=\operatorname{diag}(1,i)$ preserves the six stabiliser rays; restricting
the state set does not make that matrix nonlinear \autocite{backens2014zx}.
A restriction to standard quantum preparations, operations and measurements
cannot create a signalling protocol unavailable in the full theory. These
known quantum facts provide a counterexample to the inference, not premises
for the reconstruction.

Simon, Bužek and Gisin derive linear, completely positive density-matrix
dynamics using Hilbert-space states, the usual trace probability rule, and
no signalling \autocite{simon2001nosignaling}. Their theorem cannot be moved
upstream of the metric and Born rule without replacing those premises.
An independently formulated operational no-signalling condition does not
by itself exclude the restricted subtheory just described.

\subsubsection{Diagonalisation or inexhaustible self-ignorance supplies continuity}
\label{app:continuity-diagonal}

\noindent\textsc{Structural.}
Diagonalisation establishes failure of exhaustive representation under
specified assumptions, not connectedness \autocite{Yanofsky2003}. For example,
the space of infinite binary sequences is uncountable but totally disconnected:
two distinct sequences differ at a coordinate whose two values separate
them into disjoint open-and-closed sets. Uncountably many possibilities
need not admit a nonconstant continuous path. Density in a larger space
also does not establish that missing limits are physically admitted.

A continuum of outcome probabilities is insufficient for a different
reason. A classical bit admits every $p\in[0,1]$ through mixtures while
retaining two pure states. Likewise, analytic continuation of a logarithm
requires a specified domain and regularity; it does not establish that
physical amplitudes occupy that domain. The value-continuum premise used
in the Shannon/logarithm construction therefore cannot also count as a
derivation of connected pure states.

\subsubsection{Observer-dependent pure and mixed descriptions provide purification}
\label{app:continuity-purification}

\noindent\textsc{Structural.}
Conditioning must be distinguished from taking a subsystem. After Alice
measures her half of $(|00\rangle+|11\rangle)/\sqrt{2}$ in the computational
basis and learns the result, she assigns Bob the corresponding conditional
pure state. Without the result, Bob assigns his qubit $I/2$. A pure
description of the combined system does not make Bob's local state pure.
These standard quantum facts clarify the example; they are not pre-Born
premises for the continuity argument.

Hilbert-space purification says that every mixed state is the marginal of a
pure state of a larger system; its usual construction already uses the
Hilbert structure. Operational purification can instead be postulated before
that structure, as in Chiribella, D'Ariano and Perinotti's programme
\autocite{chiribella2010purification}. Their reconstruction also uses other
informational axioms \autocite{chiribella2011informational}. Purification has
not here been derived from the impossibility of complete self-description.

\subsubsection{No distinguished step implies arbitrary physical roots}
\label{app:continuity-roots}

\noindent\textsc{Structural.}
Under the finite-dimensional compact-group assumptions below, divisibility
of $G$ --- every element has an $n$-th root for every positive integer $n$ ---
implies connectedness. The gap is physical: absence of a distinguished step
does not establish the availability of all those roots. A formal matrix
root need not be an allowed operation. Proposition~\ref{prop:continuity-review-divisible}
retains the conditional result without supplying that premise.

\subsubsection{Connected pure states and symmetry imply continuous reversibility}
\label{app:continuity-connected}

\noindent\textsc{Structural.}
This route works conditionally. A closed group of reversible affine
transformations of a finite-dimensional compact convex state body, acting
transitively on connected pure states, has an identity component that
already acts transitively. Proposition~\ref{prop:continuity-review-connected}
gives the proof. The physical origin of connectedness remains unproved,
with closedness and transitivity also requiring explicit treatment.

Admitting all complex directions would provide a connected space even
before choosing $L^2$: the unit sphere of any norm on $\mathbb C^n$, for
$n\geq1$, is connected by radial rescaling of an ordinary sphere. A complex
representation, however, does not establish that admission. Closure under
limits does not suffice either: finite state sets and finite transformation
groups are already closed.

\subsubsection{Self-consistency requires pure fixed points}
\label{app:continuity-fixed-points}

\noindent\textsc{Structural.}
The fixed-point route proposes
\begin{equation}
  \forall g\in G\quad\exists x\in P:\ gx=x.
  \label{eq:continuity-pure-fixed-point}
\end{equation}
With compact Lie structure, transitivity and at least two pure states,
this forces a connected, positive-dimensional pure-state space
(Proposition~\ref{prop:continuity-review-fixed-points}). It is a useful
reduction of the target.

\noindent\textsc{Open.}
What remains is the derivation of \emph{pure} fixed points from physical
self-reference. A classical bit-flip has a consistent equal mixture without
a pure fixed point. Lawvere's theorem requires a suitable representation
and evaluation structure and a surjectivity hypothesis
\autocite{Lawvere_1969,Yanofsky2003}. Calling a situation self-referential
does not verify these hypotheses or put the resulting fixed point in $P$.
This route passes only if the physical and mathematical requirements are
established without assuming the full Hilbert space. It fails if consistency
supplies only mixed states or if the representation hypothesis is merely
postulated.

\subsection{The possible missing physical insight}
\label{app:continuity-candidate}

\noindent\textsc{Candidate.}
The common gap is a relation between the information an embedded observer
cannot record and the physical alternatives a theory must admit. An
obstruction to a universally correct record says what cannot be completed;
continuous reversibility asserts the existence of a family of processes.
A physical principle must explain that existence.

The complete act of self-description is a possible place to seek it. The
target, the memory describing it, the physical encoding and the feedback
that uses the description must belong to the same account. Adding memory
changes the target; redescribing the enlarged target is another physical
act. The relevant constraint may concern information, disturbance and
recoverable correlations among these parts when complete internal recording
is unavailable. It has not been shown that these correlations are phases,
that they form a continuum, or that their transformations are reversible.

One possible intermediate result would link an obstruction to continuous
reversible access to an internally available, repeatably readable distinction.
If that link followed from the observer model, the absence of the relevant
record could constrain the decomposition of pure states. The measurement
and copying operations must be derived; a partition alone does not supply
them. Another possibility is that the same analysis forces an operational
purification or reversible-dilation principle. That direction would require
an audit of the remaining reconstruction premises. Neither possibility is
presently a theorem.

\subsection{The open problem in operational form}
\label{app:continuity-open}

\noindent\textsc{Open.}
The following formulation sharpens Open Problem~\ref{op:limit-realisability}.
Work with a finite-dimensional operational theory. Declare that preparations
form a compact convex set $\Omega$, a finite collection of measurement
probabilities separates its states, and deterministic reversible operations
act affinely. Put $P=\operatorname{ext}(\Omega)$. Take complex amplitudes
only with the algebraic and compositional properties actually obtained
from the earlier self-reference argument; do not add full superposition
admissibility.

\paragraph{Physical premise.}
Specify an operational condition $\mathcal R$ that states the limits of
an observer's self-description. Include the observer, memory, record production and
feedback in the same theory. Identify the prediction tasks admitting a
diagonal response and the partial or retrospective records that remain
possible. Distinguish closure of the physical model under adding observers
from an observer's possession of an exhaustive internal representation.
Where more memory is needed, use a family of finite systems rather than a
finite register endowed with unlimited capacity.

\paragraph{Narrow mathematical task.}
Initially grant that $G$ is closed and transitive on $P$, explicitly as
additional premises. Determine whether $\mathcal R$, together with the
declared complex-amplitude structure, forces $P$ to be connected, or
equivalently under these assumptions forces $G_0$ to act transitively.
Deriving \eqref{eq:continuity-pure-fixed-point} would be one sufficient
route, but is not required.

\paragraph{Full physical task.}
Explain whether closedness, transitivity and available inverses also
follow from the observer analysis. Identify any that remain independent
postulates. A solution of the narrow task must not be presented as deriving
continuous reversibility from self-reference alone.

\paragraph{PASS: affirmative resolution.}
State $\mathcal R$ in independently meaningful operational terms and prove
\eqref{eq:continuity-accessibility} for every model of the declared premises.
Establish the available paths or sufficient connectedness result and identify
the precise premise violated by the finite counterexamples. Constructing
one connected model does not establish necessity.

\paragraph{PASS: negative resolution.}
Construct a model satisfying the same explicit premises, including the
physical observer and its prediction and record procedures, in which the
target fails. For the narrow task, retain closedness and transitivity. The
model must satisfy the substantive complex-amplitude requirements, not
merely encode classical labels with complex coordinates. Such a model would
establish that an additional physical principle is necessary.

\paragraph{FAIL as a derivation.}
Import the Born rule, a physical $L^2$ sphere, all projective rays, all
unitary transformations, continuous control or arbitrary physical roots as
an unstated premise; infer pure-state connectedness merely from continuously
variable probabilities; or assert that a toy model violates $\mathcal R$
without defining and testing it. A premise that only rephrases connectedness
also fails the explanatory aim even if the subsequent mathematics is correct.

\subsection{Conditional mathematical results}
\label{app:continuity-conditional-results}

\noindent\textsc{Structural.}
Let $\Omega$ be a compact convex set spanning a finite-dimensional affine
space and $G$ a closed subgroup of its affine automorphisms. The
automorphism group is compact: its centroid is a common fixed point, and
inscribed and circumscribed balls bound the matrices and their inverses.
The closed-subgroup theorem makes $G$ a compact matrix Lie group; its
identity component $G_0$ is path-connected and has finite index. Auxiliary
Euclidean coordinates in this argument do not identify physical
normalisation or probabilities with a quadratic norm. All three propositions
below use these hypotheses; arbitrary compact groups need not be Lie groups.

\begin{proposition}[Connected pure states]
\label{prop:continuity-review-connected}
If $G$ acts transitively on a connected pure-state space $P$, then $G_0$
acts transitively on $P$ and \eqref{eq:continuity-accessibility} holds.
\end{proposition}
\begin{proof}
The $G_0$-orbits partition $P$ into finitely many sets because $G/G_0$ is
finite and $G$ is transitive. Each orbit is compact and hence closed; in
this finite partition each is also open. Connectedness leaves one orbit.
For $x,y\in P$, choose $h\in G_0$ with $hx=y$ and a continuous path in
$G_0$ from $e$ to $h$. This is the required family of transformations.
The whole group $G$ need not be connected.
\end{proof}

\begin{proposition}[Divisibility]
\label{prop:continuity-review-divisible}
If $G$ is divisible, it is connected. If it also acts transitively on $P$,
then \eqref{eq:continuity-accessibility} holds.
\end{proposition}
\begin{proof}
The finite quotient $G/G_0$ inherits divisibility. In a finite group of
order $m$, every element satisfies $h^m=e$, so a nonidentity element has
no $m$-th root. The quotient must therefore be trivial. A connected Lie
group is path-connected, and transitivity gives the result.
\end{proof}

\begin{proposition}[Pure fixed points]
\label{prop:continuity-review-fixed-points}
Suppose $G$ acts transitively on $P$, with $|P|\geq2$, and every $g\in G$
fixes some point of $P$. Then $P$ is a connected manifold of positive
dimension and \eqref{eq:continuity-accessibility} holds.
\end{proposition}
\begin{proof}
The finite permutation group induced on the $G_0$-orbits is transitive.
If there were more than one orbit, this group would have an element moving
every orbit: its average number of fixed points is one by the orbit-counting
lemma, while its identity fixes more than one. Any lift to $G$ would fix no
point of $P$, contradicting the hypothesis. Thus $P$ is a single $G_0$-orbit
and is connected. Being a nontrivial homogeneous space of a compact Lie
group, it has positive dimension. Proposition~\ref{prop:continuity-review-connected}
supplies the continuous reversible paths.
\end{proof}

Averaging any state over a compact group gives an invariant state of
$\Omega$, which may be mixed. Consistency in the convex state space is
therefore much weaker than the pure fixed-point premise. A physical reason
to require an admissible pure fixed point is precisely what that route
still owes.
% END CONTINUITY SELF REFERENCE APPENDIX

% Main-text reference: Appendix~\ref{app:continuity-self-reference}.
% Five citation keys await import of
% bib/continuity_open_problem_refs_2026-10-18T1811.bib and BBT re-export:
% backens2014zx, bennett2003notes, simon2001nosignaling,
% chiribella2010purification, chiribella2011informational.
% The two Chiribella et al. citations support only the limited statements
% verified from their abstracts; no detailed reconstruction audit is claimed.
%
% SUGGESTED MAIN-TEXT SUMMARY (for the later main-text edit):
% Continuous reversible accessibility remains \textsc{open}. Complex
% amplitudes, probability normalisation and symmetry among admitted bases
% do not individually establish physically available paths between pure
% states. With a closed, transitive reversible group on a finite-dimensional
% compact convex state body, connectedness of the pure states would suffice;
% the outstanding task is to justify that connectedness physically.
%
% \textsc{Candidate.} The missing link may concern the complete act of
% self-description: producing a record changes the observer, and the memory,
% its encoding and the subsequent feedback must belong to the same physical
% account. A constraint on information, disturbance and recoverable
% correlations might connect the failure of complete self-recording to the
% alternatives the theory must admit. That connection is not yet a theorem.
% Appendix~\ref{app:continuity-self-reference} reviews the attempted routes,
% retains their valid conditional results, and gives explicit criteria for
% either a derivation or a countermodel. The argument must precede the Born
% rule and the physical $L^2$ metric.

\section{Continuity from self-reference: attempted derivations and the open problem}
\label{app:continuity-self-reference}

\noindent\textsc{Open.}
Can the physical limitations on an embedded observer's self-description
explain why any two pure states can be connected by a continuous family of
physically available reversible transformations? This appendix reviews the
attempts underlying Open Problem~\ref{op:limit-realisability}. Some proposed
implications fail; others give useful conditional mathematics while leaving
a physical premise unproved. The distinction matters: replacing that
premise by a mathematical reformulation does not explain why nature must
satisfy it. A status label applies to the discussion it introduces until
a different status is stated.

\subsection{The target and the order of dependence}
\label{app:continuity-target}

\noindent\textsc{Structural.}
Hardy's fifth axiom asks for continuous reversible accessibility between
any two pure states \autocite{Hardy2001a}. Let $P$ denote the admitted pure
states and $G$ the physically allowed reversible transformations. The target is
\begin{equation}
  \forall x,y\in P\quad\exists\,g:[0,1]\longrightarrow G
  \quad\text{continuous},\qquad g(0)=e,\quad g(1)x=y.
  \label{eq:continuity-accessibility}
\end{equation}
This is stronger than continuity of an individual map on states: the
intermediate transformations must also be available. The parameter need
not already be a physical time coordinate, and differentiability is not
an initial premise.

A topology can be specified before the Born rule. In a finite-dimensional
operational description, a finite list of measurement probabilities
separates states, and convergence means convergence of those probabilities.
Pure states are the extreme points of the admitted convex preparation set,
not automatically all rays of a vector space. This specifies what
connectedness would mean without assuming it \autocite{Hardy2001a}.

\noindent\textsc{Candidate.}
The earlier complex-amplitude construction is taken here as a conditional
starting point. The proposed order in \S\ref{sec:reversible-continuity-l2} is
\begin{equation}
\begin{aligned}
 &\text{self-reference and complex amplitudes}\\
 &\qquad\longrightarrow\text{continuous reversible accessibility}\\
 &\qquad\longrightarrow L^2
   \longrightarrow (S^1\longrightarrow S^3\longrightarrow S^2)
   \longrightarrow\text{Born probabilities}.
\end{aligned}
\label{eq:continuity-dependency}
\end{equation}
Each downstream arrow retains its own premises, including assumptions
about how states are represented and which transformations preserve the norm. Neither the
Born rule nor a physically normalised Hilbert sphere can therefore justify
the earlier continuity step.

\subsection{What self-recording prevents}
\label{app:continuity-self-recording}

\noindent\textsc{Structural.}
Writing a new record into an observer's memory changes its total physical
state, schematically $s\xrightarrow{\mathrm{write}\ r}s_r$. The record need
not thereby become false: it can describe a past state, a partial property,
or a predictable consequence of the writing. The stronger obstruction
concerns a universally successful predictor whose predictions can themselves
enter the behaviour predicted. Popper includes calculation and the physical
production of predictions in his analysis and establishes limitations in a
classical setting \autocite{Popper1950a}. This obstruction alone therefore
does not distinguish quantum theory from classical theories.

\noindent\textsc{Candidate.}
\textsc{Fanout} denotes the formation of an accessible, copyable outcome
record. The proposed pre-\textsc{fanout} regime lacks a committed outcome
record for the interaction in question. Turning that statement into the
absence of every preferred basis or invariant distinction requires an
additional argument. In particular, the impossibility of complete
self-recording is not the impossibility of recording any property.

\noindent\textsc{Rhyme.}
The liar and truth-teller suggest a comparison between a Boolean
self-assignment that is inconsistent and one that is consistent but
underdetermined. This
motivates the description of pre-\textsc{fanout} paradox and post-\textsc{fanout}
hypodox. It supplies neither a topology nor a probability law for the record
selected. The comparison with information in a closed timelike loop has
the same limited role.

\noindent\textsc{Structural.}
Inability to predict an outcome from the observer's information does not
establish absence of physical antecedents. A restriction on prediction
cannot alone support that causal conclusion or a continuity theorem.

\subsection{Attempted routes and their remaining gaps}
\label{app:continuity-routes}

\subsubsection{Complex amplitudes provide a continuum}
\label{app:continuity-complex}

\noindent\textsc{Structural.}
A continuous scalar field permits continuous families of vectors and
matrices; it does not make all of them physical. The complex stabiliser
fragment supplies a direct counterexample to that inference. Its
single-qubit pure states are the six rays represented by
$(1,0)$, $(0,1)$, $(1,\pm1)$ and $(1,\pm i)$, while its Clifford
transformations have 24 distinct actions on rays. It has complex amplitudes
and reversible linear maps without a continuous path between distinct pure
states \autocite{backens2014zx}.

Spekkens's toy bit is an earlier finite example, but is not itself a native
complex-amplitude model \autocite{spekkens2007toy}; the stabiliser fragment
addresses that objection. These examples test particular implications.
They have not thereby been checked against a fully specified physical model
of a self-describing observer.

\subsubsection{Probability normalisation forces unitarity and continuity}
\label{app:continuity-normalisation}

\noindent\textsc{Structural.}
For an exhaustive measurement, normalisation gives
\begin{equation}
  \sum_i p(i\mid s,M)=1.
  \label{eq:continuity-operational-normalisation}
\end{equation}
These are real outcome probabilities. Identifying them with $|a_i|^2$
already supplies the quadratic rule that the proposed order of dependence
places later. The attempted conservation argument made this identification
as a premise; its calculation is conditional on that premise and cannot
establish the earlier continuity step.

Even granting a quadratic norm and complex-linear dynamics, norm
preservation gives $U^\dagger U=I$, not the physical availability of every
element of $U(n)$. A finite subgroup also preserves the norm. Continuity of
each linear map remains distinct from a continuous path of allowed maps
connecting it to the identity. Both gaps survive the normalisation calculation.

\subsubsection{No record means no preferred basis and smooth access between bases}
\label{app:continuity-bases}

\noindent\textsc{Structural.}
Absence of a particular record does not exclude a restriction in the laws.
Symmetry among admitted bases does not imply that every mathematical basis
is admitted. Transitivity among admitted states does not imply continuous
transitivity: a finite group can move any member of a finite set to any
other. A passive change of coordinates is also distinct from a physical
operation. The stabiliser fragment treats its three Pauli bases symmetrically
while restricting the family of bases \autocite{backens2014zx}.

An invariant mathematical partition is not automatically a readable
classical record. An argument using that identification must construct an
allowed measurement that distinguishes the parts and an allowed procedure
that copies the result. Naming the parts supplies neither operation. Thus
the claim that continuity forbids a privileged finite basis does not
establish its converse.

\subsubsection{No thermodynamic cost means reversible evolution}
\label{app:continuity-thermodynamics}

\noindent\textsc{Structural.}
Logical invertibility, absence of heat dissipation and thermodynamic
reversibility are different conditions. Landauer reasoning tracks
information transferred to other degrees of freedom; it does not derive
microscopic reversible laws from zero heat. Copying classical data into a
blank register can itself be implemented reversibly \autocite{bennett2003notes}.
A cost assigned to \textsc{fanout} consequently cannot by itself prove
that all pre-\textsc{fanout} operations possess physically available
inverses, still less that those inverses belong to continuous paths.

\subsubsection{Excluded superpositions cause nonlinear dynamics and signalling}
\label{app:continuity-signalling}

\noindent\textsc{Structural.}
The first implication is false. The complex-linear reversible matrix
$S=\operatorname{diag}(1,i)$ preserves the six stabiliser rays; restricting
the state set does not make that matrix nonlinear \autocite{backens2014zx}.
A restriction to standard quantum preparations, operations and measurements
cannot create a signalling protocol unavailable in the full theory. These
known quantum facts provide a counterexample to the inference, not premises
for the reconstruction.

Simon, Bužek and Gisin derive linear, completely positive density-matrix
dynamics using Hilbert-space states, the usual trace probability rule, and
no signalling \autocite{simon2001nosignaling}. Their theorem cannot be moved
upstream of the metric and Born rule without replacing those premises.
An independently formulated operational no-signalling condition does not
by itself exclude the restricted subtheory just described.

\subsubsection{Diagonalisation or inexhaustible self-ignorance supplies continuity}
\label{app:continuity-diagonal}

\noindent\textsc{Structural.}
Diagonalisation establishes failure of exhaustive representation under
specified assumptions, not connectedness \autocite{Yanofsky2003}. For example,
the space of infinite binary sequences is uncountable but totally disconnected:
two distinct sequences differ at a coordinate whose two values separate
them into disjoint open-and-closed sets. Uncountably many possibilities
need not admit a nonconstant continuous path. Density in a larger space
also does not establish that missing limits are physically admitted.

A continuum of outcome probabilities is insufficient for a different
reason. A classical bit admits every $p\in[0,1]$ through mixtures while
retaining two pure states. Likewise, analytic continuation of a logarithm
requires a specified domain and regularity; it does not establish that
physical amplitudes occupy that domain. The value-continuum premise used
in the Shannon/logarithm construction therefore cannot also count as a
derivation of connected pure states.

\subsubsection{Observer-dependent pure and mixed descriptions provide purification}
\label{app:continuity-purification}

\noindent\textsc{Structural.}
Conditioning must be distinguished from taking a subsystem. After Alice
measures her half of $(|00\rangle+|11\rangle)/\sqrt{2}$ in the computational
basis and learns the result, she assigns Bob the corresponding conditional
pure state. Without the result, Bob assigns his qubit $I/2$. A pure
description of the combined system does not make Bob's local state pure.
These standard quantum facts clarify the example; they are not pre-Born
premises for the continuity argument.

Hilbert-space purification says that every mixed state is the marginal of a
pure state of a larger system; its usual construction already uses the
Hilbert structure. Operational purification can instead be postulated before
that structure, as in Chiribella, D'Ariano and Perinotti's programme
\autocite{chiribella2010purification}. Their reconstruction also uses other
informational axioms \autocite{chiribella2011informational}. Purification has
not here been derived from the impossibility of complete self-description.

\subsubsection{No distinguished step implies arbitrary physical roots}
\label{app:continuity-roots}

\noindent\textsc{Structural.}
Under the finite-dimensional compact-group assumptions below, divisibility
of $G$ --- every element has an $n$-th root for every positive integer $n$ ---
implies connectedness. The gap is physical: absence of a distinguished step
does not establish the availability of all those roots. A formal matrix
root need not be an allowed operation. Proposition~\ref{prop:continuity-review-divisible}
retains the conditional result without supplying that premise.

\subsubsection{Connected pure states and symmetry imply continuous reversibility}
\label{app:continuity-connected}

\noindent\textsc{Structural.}
This route works conditionally. A closed group of reversible affine
transformations of a finite-dimensional compact convex state body, acting
transitively on connected pure states, has an identity component that
already acts transitively. Proposition~\ref{prop:continuity-review-connected}
gives the proof. The physical origin of connectedness remains unproved,
with closedness and transitivity also requiring explicit treatment.

Admitting all complex directions would provide a connected space even
before choosing $L^2$: the unit sphere of any norm on $\mathbb C^n$, for
$n\geq1$, is connected by radial rescaling of an ordinary sphere. A complex
representation, however, does not establish that admission. Closure under
limits does not suffice either: finite state sets and finite transformation
groups are already closed.

\subsubsection{Self-consistency requires pure fixed points}
\label{app:continuity-fixed-points}

\noindent\textsc{Structural.}
The fixed-point route proposes
\begin{equation}
  \forall g\in G\quad\exists x\in P:\ gx=x.
  \label{eq:continuity-pure-fixed-point}
\end{equation}
With compact Lie structure, transitivity and at least two pure states,
this forces a connected, positive-dimensional pure-state space
(Proposition~\ref{prop:continuity-review-fixed-points}). It is a useful
reduction of the target.

\noindent\textsc{Open.}
What remains is the derivation of \emph{pure} fixed points from physical
self-reference. A classical bit-flip has a consistent equal mixture without
a pure fixed point. Lawvere's theorem requires a suitable representation
and evaluation structure and a surjectivity hypothesis
\autocite{Lawvere_1969,Yanofsky2003}. Calling a situation self-referential
does not verify these hypotheses or put the resulting fixed point in $P$.
This route passes only if the physical and mathematical requirements are
established without assuming the full Hilbert space. It fails if consistency
supplies only mixed states or if the representation hypothesis is merely
postulated.

\subsection{The possible missing physical insight}
\label{app:continuity-candidate}

\noindent\textsc{Candidate.}
The common gap is a relation between the information an embedded observer
cannot record and the physical alternatives a theory must admit. An
obstruction to a universally correct record says what cannot be completed;
continuous reversibility asserts the existence of a family of processes.
A physical principle must explain that existence.

The complete act of self-description is a possible place to seek it. The
target, the memory describing it, the physical encoding and the feedback
that uses the description must belong to the same account. Adding memory
changes the target; redescribing the enlarged target is another physical
act. The relevant constraint may concern information, disturbance and
recoverable correlations among these parts when complete internal recording
is unavailable. It has not been shown that these correlations are phases,
that they form a continuum, or that their transformations are reversible.

One possible intermediate result would link an obstruction to continuous
reversible access to an internally available, repeatably readable distinction.
If that link followed from the observer model, the absence of the relevant
record could constrain the decomposition of pure states. The measurement
and copying operations must be derived; a partition alone does not supply
them. Another possibility is that the same analysis forces an operational
purification or reversible-dilation principle. That direction would require
an audit of the remaining reconstruction premises. Neither possibility is
presently a theorem.

\subsection{The open problem in operational form}
\label{app:continuity-open}

\noindent\textsc{Open.}
The following formulation sharpens Open Problem~\ref{op:limit-realisability}.
Work with a finite-dimensional operational theory. Declare that preparations
form a compact convex set $\Omega$, a finite collection of measurement
probabilities separates its states, and deterministic reversible operations
act affinely. Put $P=\operatorname{ext}(\Omega)$. Take complex amplitudes
only with the algebraic and compositional properties actually obtained
from the earlier self-reference argument; do not add full superposition
admissibility.

\paragraph{Physical premise.}
Specify an operational condition $\mathcal R$ that states the limits of
an observer's self-description. Include the observer, memory, record production and
feedback in the same theory. Identify the prediction tasks admitting a
diagonal response and the partial or retrospective records that remain
possible. Distinguish closure of the physical model under adding observers
from an observer's possession of an exhaustive internal representation.
Where more memory is needed, use a family of finite systems rather than a
finite register endowed with unlimited capacity.

\paragraph{Narrow mathematical task.}
Initially grant that $G$ is closed and transitive on $P$, explicitly as
additional premises. Determine whether $\mathcal R$, together with the
declared complex-amplitude structure, forces $P$ to be connected, or
equivalently under these assumptions forces $G_0$ to act transitively.
Deriving \eqref{eq:continuity-pure-fixed-point} would be one sufficient
route, but is not required.

\paragraph{Full physical task.}
Explain whether closedness, transitivity and available inverses also
follow from the observer analysis. Identify any that remain independent
postulates. A solution of the narrow task must not be presented as deriving
continuous reversibility from self-reference alone.

\paragraph{PASS: affirmative resolution.}
State $\mathcal R$ in independently meaningful operational terms and prove
\eqref{eq:continuity-accessibility} for every model of the declared premises.
Establish the available paths or sufficient connectedness result and identify
the precise premise violated by the finite counterexamples. Constructing
one connected model does not establish necessity.

\paragraph{PASS: negative resolution.}
Construct a model satisfying the same explicit premises, including the
physical observer and its prediction and record procedures, in which the
target fails. For the narrow task, retain closedness and transitivity. The
model must satisfy the substantive complex-amplitude requirements, not
merely encode classical labels with complex coordinates. Such a model would
establish that an additional physical principle is necessary.

\paragraph{FAIL as a derivation.}
Import the Born rule, a physical $L^2$ sphere, all projective rays, all
unitary transformations, continuous control or arbitrary physical roots as
an unstated premise; infer pure-state connectedness merely from continuously
variable probabilities; or assert that a toy model violates $\mathcal R$
without defining and testing it. A premise that only rephrases connectedness
also fails the explanatory aim even if the subsequent mathematics is correct.

\subsection{Conditional mathematical results}
\label{app:continuity-conditional-results}

\noindent\textsc{Structural.}
Let $\Omega$ be a compact convex set spanning a finite-dimensional affine
space and $G$ a closed subgroup of its affine automorphisms. The
automorphism group is compact: its centroid is a common fixed point, and
inscribed and circumscribed balls bound the matrices and their inverses.
The closed-subgroup theorem makes $G$ a compact matrix Lie group; its
identity component $G_0$ is path-connected and has finite index. Auxiliary
Euclidean coordinates in this argument do not identify physical
normalisation or probabilities with a quadratic norm. All three propositions
below use these hypotheses; arbitrary compact groups need not be Lie groups.

\begin{proposition}[Connected pure states]
\label{prop:continuity-review-connected}
If $G$ acts transitively on a connected pure-state space $P$, then $G_0$
acts transitively on $P$ and \eqref{eq:continuity-accessibility} holds.
\end{proposition}
\begin{proof}
The $G_0$-orbits partition $P$ into finitely many sets because $G/G_0$ is
finite and $G$ is transitive. Each orbit is compact and hence closed; in
this finite partition each is also open. Connectedness leaves one orbit.
For $x,y\in P$, choose $h\in G_0$ with $hx=y$ and a continuous path in
$G_0$ from $e$ to $h$. This is the required family of transformations.
The whole group $G$ need not be connected.
\end{proof}

\begin{proposition}[Divisibility]
\label{prop:continuity-review-divisible}
If $G$ is divisible, it is connected. If it also acts transitively on $P$,
then \eqref{eq:continuity-accessibility} holds.
\end{proposition}
\begin{proof}
The finite quotient $G/G_0$ inherits divisibility. In a finite group of
order $m$, every element satisfies $h^m=e$, so a nonidentity element has
no $m$-th root. The quotient must therefore be trivial. A connected Lie
group is path-connected, and transitivity gives the result.
\end{proof}

\begin{proposition}[Pure fixed points]
\label{prop:continuity-review-fixed-points}
Suppose $G$ acts transitively on $P$, with $|P|\geq2$, and every $g\in G$
fixes some point of $P$. Then $P$ is a connected manifold of positive
dimension and \eqref{eq:continuity-accessibility} holds.
\end{proposition}
\begin{proof}
The finite permutation group induced on the $G_0$-orbits is transitive.
If there were more than one orbit, this group would have an element moving
every orbit: its average number of fixed points is one by the orbit-counting
lemma, while its identity fixes more than one. Any lift to $G$ would fix no
point of $P$, contradicting the hypothesis. Thus $P$ is a single $G_0$-orbit
and is connected. Being a nontrivial homogeneous space of a compact Lie
group, it has positive dimension. Proposition~\ref{prop:continuity-review-connected}
supplies the continuous reversible paths.
\end{proof}

Averaging any state over a compact group gives an invariant state of
$\Omega$, which may be mixed. Consistency in the convex state space is
therefore much weaker than the pure fixed-point premise. A physical reason
to require an admissible pure fixed point is precisely what that route
still owes.
% END CONTINUITY SELF REFERENCE APPENDIX
